\begin{definition} \textbf{Conjunto acotado.} \label{def:conj_acotado}
    Sea $A \subseteq \Rn{n}$. Decimos que $A$ es un \emph{conjunto acotado} si existe $\delta > 0$ tal que la bola de radio $\delta$ centrada en el origen contiene al conjunto $A$, es decir, si:
    \[
     \exists \delta > 0 \; \colon \; A \subseteq B_{\delta}(\mathbf{0}).
    \]
    % Equivalentemente, $A$ es un \emph{conjunto acotado} si $\exists \delta > 0, \vect{x}_0 \in \Rn{n} \; \colon \; A \subseteq B_{\delta}(\vect{x}_0)$.
\end{definition}
  



